\subsection{Inference a parametry generování}
\label{subsec:inference-generation}

Při inferenci model generuje výstup autoregresivně, v každém kroku predikuje pravděpodobnostní rozdělení přes celý slovník tokenů a vybere z něj jeden token, který přidá k dosavadní sekvenci.
Způsob výběru tokenu určuje \textit{dekódovací strategie}:

\begin{itemize}
    \item \textbf{Greedy decoding} -- vždy vybere nejpravděpodobnější token.
    Což je deterministické, avšak náchylné k opakování.

    \item \textbf{Vzorkování s teplotou} -- vybírá token náhodně z rozdělení upraveného parametrem teploty $T$.
    Nízká teplota ($T < 1$) zvyšuje pravděpodobnost nejpravděpodobnějšího tokenu, vysoká teplota ($T > 1$) rozdělení uhlazuje a zvyšuje variabilitu.

    \item \textbf{Top-p (nucleus) sampling} -- vybírá token z nejmenší podmnožiny tokenů, jejichž kumulativní pravděpodobnost překračuje práh $p$, čímž zabraňuje výběru velmi nepravděpodobných tokenů při zachování jisté míry variability.
\end{itemize}

Pro analýzu zdrojového kódu ve školním prostředí je žádoucí konzistentnost a determinismus, kde cílem je diagnostika, nikoli kreativní generování.
Proto se doporučuje nízká teplota v rozsahu $T = 0{,}1$ až $T = 0{,}3$ kombinovaná s top-p vzorkováním.

Důležitou optimalizací inference je \textbf{KV cache}: klíče a hodnoty ($\mathbf{K}$, $\mathbf{V}$) pro tokeny vstupního prefixu jsou uloženy do cache a nepřepočítávají se při každém dalším generovacím kroku.
Tím se drasticky snižuje výpočetní náročnost inference, avšak za cenu vyšší spotřeby paměti.
KV cache roste lineárně s délkou sekvence a počtem vrstev modelu.

\subsubsection*{Zjednodušený přehled procesu}
\label{subsubsec:llm-simplified-process}

Prakticky lze z hlediska fungování modelu proces zjednodušit do následujících kroků:

\begin{enumerate}
    \item vstupní text je rozdělen na tokeny,
    \item jednotlivé tokeny jsou převedeny na numerické reprezentace (embeddingy) včetně poziční informace,
    \item data procházejí opakovaně několika bloky architektury Transformer (self-attention a feed-forward vrstvy),
    \item výstupní vrstva vypočítá pravděpodobnostní rozdělení nad možnými dalšími tokeny,
    \item na základě zvoleného dekódovacího mechanismu je vybrán konkrétní token,
    \item proces se opakuje až do splnění ukončující podmínky (např.\ vygenerování speciálního koncového tokenu nebo dosažení maximální délky sekvence).
\end{enumerate}

\subsubsection*{Rozdělení podle způsobu trénování}
\label{subsubsec:llm-training-types}

Z hlediska trénování lze modely rozdělit do několika skupin:

\begin{itemize}
    \item \textbf{Předtrénované modely} -- univerzální modely trénované na rozsáhlých datech, které lze použít pro široké spektrum úloh bez dalšího trénování.
    \item \textbf{Jemně doladěné (fine-tuned) modely} -- modely, které jsou dále specializovány na konkrétní úlohu nebo doménu pomocí dodatečných trénovaných dat.
    \item \textbf{Multimodální modely} -- modely schopné pracovat s více typy vstupních dat (např.\ text a obraz).
    \item \textbf{Doménově specifické modely} -- modely optimalizované pro konkrétní průmyslové nebo odborné oblasti.
\end{itemize}

Z pohledu této diplomové práce jsou nejrelevantnější autoregresivní transformer modely, které umožňují generovat textové výstupy na základě vstupního zdrojového kódu.
Tyto modely jsou vhodné pro generování komentářů, shrnutí řešení a identifikaci potenciálních problémů ve studentských úlohách.

\endinput